%% Full selected Section 3 proofs, reconstructed from arXiv:0808.0521v1.
\begin{annexitem}{annex:p2.syllogistic.identifications}{thm:p02-syllogistic-semantic-identifications}{Complements, formula identifications and first-order translation}
\begin{proof}
For every literal \(l\), the interpretation of \(\bar l\) is
\(A\setminus l^{\mathfrak A}\), including the case in which \(l\) is already
negative. Hence \(l^{\mathfrak A}\not\subseteq m^{\mathfrak A}\) holds
exactly when \(l^{\mathfrak A}\cap(\bar m)^{\mathfrak A}\ne\varnothing\).
This proves the formula-complement assertion for universal sentences.
For existential sentences, the intersection is empty exactly when
\(l^{\mathfrak A}\subseteq(\bar m)^{\mathfrak A}\), proving the other case.
Applying the formula bar twice restores the sentence.

Intersection is symmetric. Also, for subsets \(L,M\subseteq A\),
\(L\subseteq M\) if and only if \(A\setminus M\subseteq A\setminus L\).
These facts prove the two formula identifications. An existential sentence
lies in \(\mathcal S\) precisely when at least one of its two literals is
positive. A universal sentence lies in \(\mathcal S\) precisely when its
first literal is positive or its second literal is negative. The respective
identifications preserve these conditions. All such sentences are permitted
in \(\mathcal S^\dagger\).

The bar of an existential representative and the bar of its swapped
representative are universal contrapositives. The bars of universal
contrapositives are swapped existential representatives. Thus formula
complementation is well defined under the identifications. Fragment closure
under complementation follows directly from the four forms in
Definition~\ref{def:p02-syllogistic-language}.

No element belongs to both \(l^{\mathfrak A}\) and its complement, so an
absurdity is always false. Finally, interpret the first-order predicate
\(p(x)\) by membership in \(p^{\mathfrak A}\). The universal translation
says exactly that every member of the first literal's extension belongs to
the second. The existential translation says exactly that the intersection
has a member. This also proves the translation claim when either literal is
negative.
\end{proof}
\end{annexitem}

\begin{annexitem}{annex:p2.syllogistic.rules}{thm:p02-syllogistic-rules}{Every rule case, restriction and the redundancy of D1}
\begin{proof}
We first check soundness in a fixed structure. In the next two paragraphs,
set expressions written with literal names abbreviate their interpretations;
for example, \(l\cap n\) abbreviates
\(l^{\mathfrak A}\cap n^{\mathfrak A}\).
For \(\mathsf S^\dagger\), D preserves the witness:
an element in \(l\cap n\) also belongs to \(m\) when \(l\subseteq m\),
and therefore lies in \(m\cap n\). B is transitivity of inclusion. In A,
\(l\subseteq\bar l\) implies \(l=\varnothing\), and the empty set is
included in every \(m\). T is reflexivity of inclusion. In I, a member of
\(l\cap m\) is a member of \(l\cap l\). No structure satisfies both
\(\psi\) and \(\bar\psi\), so X is sound. In N, the inclusion
\(\bar l\subseteq l\) makes \(\bar l\) empty, hence \(l=A\). Nonemptiness
of \(A\) supplies the requested witness.

For \(\mathsf S\), B, A, T, I and X are the same set-theoretic arguments
with restricted metavariables. D1 transfers a witness in \(p\cap q\)
through \(q\subseteq l\). D2 transfers a witness in \(p\cap l\)
through \(p\subseteq q\). For D3, a witness in \(p\cap l\) cannot be a
\(q\), since \(q\subseteq\bar l\); it therefore belongs to
\(p\cap\bar q\). Atom substitution merely changes the assigned extensions,
so all substitution instances remain sound. Induction on a finite derivation
tree proves soundness of each derivability relation.

Return now to literal syntax. We check every possible
\(\mathsf S^\dagger\) rule whose formulas lie
in \(\mathcal S\). Throughout this check, the only forbidden forms are a
universal negative-to-positive inclusion and an existential pair of negative
literals.

\paragraph{D.}
Consider premises \(\exists(l,n),\forall(l,m)\) and conclusion
\(\exists(m,n)\). If \(m\) is positive, \(l\) must be positive because
\(\forall(l,m)\in\mathcal S\). This is D2. If \(m\) is negative, \(n\)
must be positive because the conclusion is in \(\mathcal S\). Replace the
universal premise by \(\forall(\bar m,\bar l)\), the existential premise
by \(\exists(n,l)\), and the conclusion by \(\exists(n,m)\). With atoms
\(n,\bar m\), this is precisely D3.

\paragraph{B.}
Consider \(\forall(l,m),\forall(m,n)\) and \(\forall(l,n)\).
If \(l,m\) are positive, the instance is B of \(\mathsf S\).
If \(l\) is negative, legality of the two premises forces \(m,n\) both
negative. Reverse the chain by contraposition to obtain
\(\forall(\bar n,\bar m),\forall(\bar m,\bar l)\), with the first two
terms positive. Its B conclusion \(\forall(\bar n,\bar l)\) is the
contrapositive of the desired sentence. The remaining case has \(l\)
positive and \(m\) negative. Then \(n\) is negative; the same reversed
chain has \(\bar n,\bar m\) positive and \(\bar l\) a permitted negative
last literal, so again it is an \(\mathsf S\) instance.

\paragraph{A, T and I.}
An A premise \(\forall(l,\bar l)\) is in \(\mathcal S\) only when \(l\)
is positive, giving its restricted A rule. A T instance with \(l=p\) is
already restricted T; with \(l=\bar p\), its conclusion is identified with
\(\forall(p,p)\). An I conclusion \(\exists(l,l)\) is in \(\mathcal S\)
only when \(l\) is positive, giving restricted I.

\paragraph{X and N.}
If all formulas in an X instance belong to \(\mathcal S\), it is exactly
restricted X. If \(l\) is positive, the N premise
\(\forall(\bar l,l)\) is outside \(\mathcal S\). If \(l\) is negative,
its conclusion \(\exists(l,l)\) is outside. There is no N instance with
every formula in \(\mathcal S\).

Conversely, each of restricted B, A, T, I and X is the corresponding
dagger rule. D2 is D. For D1, swap \(\exists(p,q)\) to
\(\exists(q,p)\), apply D with \(\forall(q,l)\), and swap the conclusion.
For D3, contrapose its universal premise to \(\forall(l,\bar q)\), swap
\(\exists(p,l)\) to \(\exists(l,p)\), apply D, and swap back.
This proves both directions of the rule-instance assertion.

Finally, D1 is redundant even in the restricted calculus. If its last literal
is positive, say \(l=o\), existential symmetry makes it a D2 instance.
If \(l=\bar o\), rewrite \(\forall(q,\bar o)\) as
\(\forall(o,\bar q)\); together with \(\exists(p,q)\), this is D3 with
conclusion \(\exists(p,\bar o)\). Both polarity cases have been covered.
\end{proof}
\end{annexitem}

\begin{annexitem}{annex:p2.syllogistic.separation}{thm:p02-syllogistic-separation}{Separation of universal constraints and witnesses}
\begin{proof}
In a disjoint union, an element of a component satisfies exactly the literals
it satisfied in that component. Indeed, its membership in a positive atom is
unchanged, and relative complementation gives the same assertion for negative
literals. Therefore a universal sentence holds in the union exactly when it
holds in every component; an existential holds exactly when some component
has a witness. In a nonempty restriction, the interpretation of each literal
is its original interpretation intersected with the smaller domain, so
universal sentences are preserved by restriction.

For part 1, suppose \(\Phi\not\models\theta\). Choose a model
\(\mathfrak A\) of \(\Phi\cup\Psi\) and a model \(\mathfrak B\) of
\(\Phi\cup\{\bar\theta\}\). Rename elements to make the domains disjoint.
The union satisfies all universals in \(\Phi\), has all witnesses required
by \(\Psi\) in its \(\mathfrak A\) component, and has a witness for the
existential \(\bar\theta\) in its \(\mathfrak B\) component. This
contradicts \(\Phi\cup\Psi\models\theta\).

For part 2, suppose no such \(\psi\) exists. For each \(\psi\in\Psi\)
choose a structure satisfying \(\Phi\cup\{\psi,\bar\theta\}\), and take
tagged copies with disjoint domains. Their union has a witness for every
member of \(\Psi\). It also satisfies all the universals in \(\Phi\)
and the universal \(\bar\theta\), because every component does. The union
is a nonempty structure since \(\Psi\ne\varnothing\). This is again a
countermodel to the assumed entailment. This argument applies to arbitrary
sets of premises in the ambient set-theoretic metatheory.

For part 3, if \(\Phi\not\models\forall(\bar l,l)\), choose a model of
\(\Phi\) containing an element \(a\) in \(\bar l\) but not in \(l\).
Its restriction to \(\{a\}\) still satisfies \(\Phi\), and no element
there belongs to \(l\cap m\). It falsifies \(\theta\). Thus
\(\Phi\models\forall(\bar l,l)\). Repeating the restriction argument
with an element outside \(m\) proves
\(\Phi\models\forall(\bar m,m)\). If \(\Phi\) has no model, both
universal consequences hold vacuously, so that case is included.
\end{proof}
\end{annexitem}

\begin{annexitem}{annex:p2.syllogistic.closure}{thm:p02-syllogistic-closure}{Closure, unions and a two-literal conflict}
\begin{proof}
For \(l\in V\), rule T derives \(\forall(l,l)\), so \(l\in V^\Phi\).
If \(n\in(V^\Phi)^\Phi\), there is \(m\in V^\Phi\) with
\(\Phi\vdash_{\mathrm{BTA}}\forall(m,n)\). Some \(l\in V\) also
satisfies \(\Phi\vdash_{\mathrm{BTA}}\forall(l,m)\). Combining the two
finite derivations by B gives \(\forall(l,n)\), and hence \(n\in V^\Phi\).
This proves closure. It also shows monotonicity: if \(V\subseteq W\), every
literal reached from a member of \(V\) is reached from a member of \(W\).

If \(m\in(\bigcup_i V_i)^\Phi\), its witnessing starting literal belongs
to some \(V_i\), so \(m\in V_i^\Phi\subseteq V_i\) whenever each \(V_i\)
is closed. Thus their union is closed; for an empty family both sides are
empty by the definition of closure.

If \(V^\Phi\) is inconsistent, take \(m,\bar m\in V^\Phi\). There are
\(l,l'\in V\) with BTA derivations of \(\forall(l,m)\) and
\(\forall(l',\bar m)\). The second sentence is identified with
\(\forall(m,\bar l')\). B yields \(\forall(l,\bar l')\), as required.
This last relation is symmetric in \(l,l'\) under universal contraposition.
\end{proof}
\end{annexitem}

\begin{annexitem}{annex:p2.syllogistic.extension}{thm:p02-syllogistic-extension}{Consistent literal completion and its singleton model}
\begin{proof}
Let \(V_0\) be nonempty, consistent and \(\Phi\)-closed. Enumerate all
literals as \(l_1,l_2,\ldots\). Given \(V_i\), leave it unchanged if
\(l_{i+1}\in V_i\); otherwise put
\(V_{i+1}=(V_i\cup\{\bar l_{i+1}\})^\Phi\).
Closure and inclusion from the preceding lemma show that this sequence is
increasing, every member is closed, and every member remains nonempty.

Suppose a first inconsistent stage occurs at \(V_{i+1}\), and write
\(r=l_{i+1}\). That stage must have been added because \(r\notin V_i\).
The conflict lemma supplies \(u,v\in V_i\cup\{\bar r\}\) with
\(\Phi\vdash_{\mathrm{BTA}}\forall(u,\bar v)\). Both cannot belong
to \(V_i\): closure would then put \(\bar v\) there together with \(v\).
By contraposition we may exchange the pair and assume \(v=\bar r\).
Thus \(\Phi\vdash_{\mathrm{BTA}}\forall(u,r)\).

If \(u\in V_i\), closure gives \(r\in V_i\), a contradiction. The only
remaining case is \(u=\bar r\), giving \(\forall(\bar r,r)\).
Choose \(w\in V_i\), using nonemptiness. Rule A, with its subject literal
\(\bar r\), derives \(\forall(\bar r,\bar w)\); universal contraposition
identifies this with \(\forall(w,r)\). Closure again gives \(r\in V_i\).
Every possible conflict is impossible, so all stages are consistent.

Let \(V=\bigcup_i V_i\). It is closed by the union assertion. If both a
literal and its complement belonged to \(V\), their occurrences would lie
together in some stage, contradicting that stage's consistency. For every
literal \(l_j\), either it was present at stage \(j-1\), or its complement
was added at stage \(j\). Therefore \(V\) is complete and contains \(V_0\).

Define a structure on \(\{a\}\) by interpreting \(p\) as \(\{a\}\) when
\(p\in V\), and as the empty set otherwise. Completeness and consistency
give exactly one of \(p,\bar p\) in \(V\), so for every literal \(l\),
\(a\in l^{\mathfrak A}\) if and only if \(l\in V\). If a premise
\(\forall(l,m)\in\Phi\) has \(l\in V\), its one-line premise derivation
and closure give \(m\in V\). If \(l\notin V\), the inclusion is vacuous.
Every universal premise is satisfied. This proves the singleton-model claim.
\end{proof}
\end{annexitem}

\begin{annexitem}{annex:p2.syllogistic.sdagger}{thm:p02-syllogistic-sdagger-completeness}{Complete direct derivations in the noun-negation fragment}
\begin{proof}
Soundness was proved for all rules and derivation trees in
Lemma~\ref{thm:p02-syllogistic-rules}. For completeness, suppose
\(\Theta\models\theta\). Translate into unary first-order logic as in
Lemma~\ref{thm:p02-syllogistic-semantic-identifications}. First-order
compactness (Theorem~\ref{thm:fol-compactness}) applied to
\(\Theta\cup\{\bar\theta\}\) provides a finite
\(\Theta_0\subseteq\Theta\) with \(\Theta_0\models\theta\).
A derivation from that subset is a derivation from \(\Theta\), so assume
from now on that \(\Theta\) is finite. Write
\(\Theta=\Phi\cup\Psi\) for its universal and existential sentences.

We will use the following immediate BTA consequence with its precise rule
steps. From \(\forall(\bar r,r)\), A gives
\(\forall(\bar r,\bar u)\) for any literal \(u\); contraposition gives
\(\forall(u,r)\). Thus a universally full literal accepts inclusion from
every literal. This is an application of A and the declared identification,
not an additional primitive rule.

First assume \(\Theta\) is satisfiable.

\paragraph{A universal conclusion.}
Write \(\theta=\forall(l_0,m_0)\). The separation lemma gives
\(\Phi\models\theta\). The set
\(\{l_0,\bar m_0\}^\Phi\) is nonempty and closed. If consistent,
the extension lemma supplies a singleton model of \(\Phi\) whose element
belongs to \(l_0\) and \(\bar m_0\); it falsifies \(\theta\).
The closure is therefore inconsistent. The conflict lemma gives
\(u,v\in\{l_0,\bar m_0\}\) with
\(\Phi\vdash_{\mathrm{BTA}}\forall(u,\bar v)\).
There are four ordered choices. The choice \(u=l_0,v=\bar m_0\) is the
required inclusion itself; the reversed mixed choice is its contrapositive.
If both are \(l_0\), we have \(\forall(l_0,\bar l_0)\), and A directly
gives \(\forall(l_0,m_0)\). If both are \(\bar m_0\), we have
\(\forall(\bar m_0,m_0)\), and the fullness argument above gives
\(\forall(l_0,m_0)\). Thus every universal consequence is derivable using
BTA alone.

\paragraph{An existential conclusion with an existential premise.}
Suppose \(\theta=\exists(l,m)\) and \(\Psi\ne\varnothing\).
By separation choose \(\psi=\exists(a,b)\in\Psi\) such that
\(\Phi\cup\{\psi\}\models\exists(l,m)\), where here \(a,b\)
are literals, not elements of a structure.

Consider \(\{a,b,\bar l\}^\Phi\). If it were consistent, its complete
extension would produce a singleton model of \(\Phi\) whose element
satisfies \(a,b,\bar l\). That model satisfies \(\psi\) and falsifies
\(\exists(l,m)\), a contradiction. Its closure is inconsistent, so
some \(u,v\in\{a,b,\bar l\}\) satisfy
\(\Phi\vdash_{\mathrm{BTA}}\forall(u,\bar v)\).
They cannot both belong to \(\{a,b\}\): in a model of the satisfiable
set \(\Theta\), the witness to \(\exists(a,b)\) would then satisfy
both \(v\) and \(\bar v\), by soundness of BTA. Interchanging the pair
by contraposition if necessary, take \(v=\bar l\).

If \(u=a\), this gives \(\Phi\vdash_{\mathrm{BTA}}\forall(a,l)\).
If \(u=b\), it gives \(\Phi\vdash_{\mathrm{BTA}}\forall(b,l)\).
If \(u=\bar l\), the fullness argument gives
\(\Phi\vdash_{\mathrm{BTA}}\forall(a,l)\). We have therefore proved
that at least one of the witness's two literals universally implies \(l\).
Repeating the complete argument with \(m\) in place of \(l\) proves that
at least one implies \(m\).

By swapping \(a,b\) in \(\psi\) when needed, arrange that
\(\Phi\vdash\forall(a,l)\). There are two cases left for \(m\).
If \(\Phi\vdash\forall(b,m)\), D applied to
\(\exists(a,b)\) and \(\forall(a,l)\) gives \(\exists(l,b)\).
Swap it to \(\exists(b,l)\), apply D with \(\forall(b,m)\), and swap
the resulting \(\exists(m,l)\) to obtain the desired sentence.
If instead \(\Phi\vdash\forall(a,m)\), first use I on \(\psi\)
to get \(\exists(a,a)\). D with \(\forall(a,l)\) gives
\(\exists(l,a)\). Swap, apply D with \(\forall(a,m)\), and swap again.
This covers both target literals arising from the same witness property
as well as their arising from different properties.

\paragraph{An existential conclusion from universals alone.}
Now let \(\Psi=\varnothing\) and \(\theta=\exists(l,m)\).
The separation lemma gives semantic consequences
\(\forall(\bar l,l)\) and \(\forall(\bar m,m)\) of \(\Phi\).
The universal case already proved derives both. N applied to the first
gives \(\exists(l,l)\). The fullness argument applied to the second
gives \(\forall(l,m)\). D and existential symmetry yield
\(\exists(l,m)\). This step uses the nonempty-universe assumption exactly
where rule N requires it.

\paragraph{Inconsistent premises.}
The preceding cases cover every conclusion from a satisfiable finite set.
If \(\Theta\) is unsatisfiable, choose an inclusion-minimal unsatisfiable
subset \(\Delta\subseteq\Theta\). It is nonempty because the empty
theory has a one-element model. Choose \(\rho\in\Delta\) and put
\(\Delta'=\Delta\setminus\{\rho\}\). Minimality makes \(\Delta'\)
satisfiable; unsatisfiability of \(\Delta'\cup\{\rho\}\) gives
\(\Delta'\models\bar\rho\). The already proved satisfiable case gives
\(\Delta'\vdash_{\mathsf S^\dagger}\bar\rho\). With the available
premise \(\rho\), X concludes \(\theta\). All premises of this finite
direct proof belong to \(\Theta\). No reductio step occurs in the proof
system: the model arguments here establish the metatheorem.
\end{proof}
\end{annexitem}

\begin{annexitem}{annex:p2.syllogistic.s}{thm:p02-syllogistic-s-completeness}{Conservativity and completeness of the smaller syllogistic}
\begin{proof}
Restricted soundness and embedding of its rules into the dagger system
follow from Lemma~\ref{thm:p02-syllogistic-rules}. It remains to convert
a semantic consequence between \(\mathcal S\)-sentences into a restricted
direct derivation. First-order compactness reduces to a finite premise
subset exactly as in the preceding proof.

Suppose first that \(\Theta\subseteq\mathcal S\) is satisfiable,
\(\theta\in\mathcal S\), and \(\Theta\models\theta\).
Dagger completeness gives a finite
\(\mathsf S^\dagger\) derivation of \(\theta\) from \(\Theta\).
For any positive atoms \(p,q\), the set \(\Theta\) cannot entail
\(\forall(\bar p,q)\). Otherwise the separation lemma would make this
a consequence of the universal subset \(\Phi\subseteq\Theta\).
Take a singleton structure with every positive atom empty. Every universal
\(\mathcal S\)-sentence is true: an inclusion with a positive subject is
vacuous, and an inclusion between two negative literals relates the entire
universe to itself. Yet \(\forall(\bar p,q)\) is false. By dagger
soundness, no such negative-to-positive universal can occur in a derivation
from \(\Theta\).

An existential sentence outside \(\mathcal S\) has two negative literals,
say \(\exists(\bar p,\bar q)\). Follow its occurrence towards the root
of a derivation tree. X cannot appear anywhere, because its contradictory
antecedents would both be true in a model of the satisfiable \(\Theta\).
Among the remaining rules, only I or D can consume this existential.
I produces another existential with two negative literals. A D use can
swap the two antecedent literals, but in either orientation it transfers
one negative literal through an inclusion with that negative subject.
The inclusion's consequent must also be negative, since negative-to-positive
universals have just been excluded. Its existential conclusion therefore
again has two negative literals. Such an occurrence propagates along every
ancestor step to the root. It cannot occur in a proof whose root is
\(\theta\in\mathcal S\).

The two excluded forms exhaust all dagger sentences outside \(\mathcal S\).
Hence every sentence in this particular proof tree belongs to
\(\mathcal S\). The complete rule-instance check in
Lemma~\ref{thm:p02-syllogistic-rules} replaces every node by a restricted
rule instance, yielding \(\Theta\vdash_{\mathsf S}\theta\).

If finite \(\Theta\) is unsatisfiable, choose a minimal unsatisfiable
\(\Delta=\Delta'\cup\{\rho\}\subseteq\Theta\), as before.
The set \(\Delta'\) is satisfiable and entails \(\bar\rho\).
The complement belongs to \(\mathcal S\), so the restricted satisfiable
case gives \(\Delta'\vdash_{\mathsf S}\bar\rho\). X with \(\rho\)
derives the desired \(\theta\). Enlarging the premise set to the original
possibly infinite \(\Theta\) preserves this finite derivation.

We have proved semantic completeness of \(\mathsf S\). Finally,
\(\Theta\vdash_{\mathsf S^\dagger}\theta\) implies
\(\Theta\models\theta\) by dagger soundness, hence
\(\Theta\vdash_{\mathsf S}\theta\) by the result just proved.
Together with the rule embedding this establishes conservativity, including
the case of inconsistent premises.
\end{proof}
\end{annexitem}

\begin{annexitem}{annex:p2.syllogistic.saturation}{thm:p02-syllogistic-saturation}{Finite vocabulary, saturation and polynomial time}
\begin{proof}
An input conclusion contains an atom, so \(Q\ne\varnothing\). Choose
\(p_0\in Q\). In any finite derivation of the input conclusion, replace
every atom outside \(Q\) by \(p_0\), while fixing all atoms in \(Q\).
This is an allowed atom substitution. It preserves every rule instance and
every formula identification, and it leaves the input premises and conclusion
unchanged. Distinctness of atoms is not a side condition in either sheet.
Thus a derivable input conclusion has a proof using only \(Q\).

Writing \(n=|Q|\), there are \(2n\) literals and at most
\(2(2n)^2=8n^2\) ordered sentences before identifications. The smaller
fragment is a subset. Canonical representatives for symmetry and
contraposition can be chosen by a fixed lexicographic order on the input
atom indices and the finitely many syntactic forms.

All sheets are fixed and each of their rule patterns uses a fixed number of
literal or formula positions. Enumerating their instances over the at most
\(8n^2\) sentences therefore takes polynomial time. For example, the largest
unoptimized enumeration for X considers a choice of \(\psi\) and conclusion
\(\varphi\); even enumerating all three formula positions independently
would still be polynomial. Testing identification, complementation and
antecedent membership is polynomial in the finite input encoding.

At most \(8n^2\) successful additions can occur. An iteration over all
instances until closure thus terminates in polynomial time. Induction on
successful additions supplies a direct proof from the stored predecessor
records. Conversely, induction on a finite proof over \(Q\) shows that
every node enters the saturated set: leaves enter initially or by their
zero-premise rule, and any internal node enters once its antecedents do.
Consequently saturation contains \(\theta\) exactly when it is derivable.
The two completeness theorems equate this with semantic entailment.

If \(\theta\) is absent, semantic non-entailment means there exists a
structure satisfying the premises and falsifying it. The procedure as stated
decides this fact; it does not claim to output that structure unless a separate
countermodel extraction is performed. If desired, a finite search suffices:
any model can be replaced by one representative of each realized truth pattern
on \(Q\), preserving all relevant unary sentences, so at most \(2^n\)
elements need be considered. This last observation is a model-size bound,
not the polynomial saturation algorithm's running-time claim.
\end{proof}
\end{annexitem}
